% =====================================================================
%  1 Introdução
% =====================================================================
\section{Introdução}
\label{sec:introducao}

Assistentes de voz instalados em alto-falantes inteligentes, como a Amazon
Alexa nos dispositivos Echo, passaram a fazer parte do cotidiano doméstico.
Estudos baseados em registros de uso mostram que ouvir música, fazer buscas
e controlar outros dispositivos da casa estão entre as atividades mais
frequentes com esses assistentes~\cite{ammari2019music}. Diferentemente de
um smartphone, o alto-falante costuma ficar em um espaço comum da casa, e
seu uso se entrelaça com as conversas e as atividades simultâneas dos
moradores, como as refeições em família~\cite{porcheron2018voice}.

Entre os recursos oferecidos pela Alexa estão as Rotinas, que permitem
associar um evento disparador (um horário, um comando de voz, um alarme) a
uma sequência de ações executadas pelo assistente e pelos dispositivos
conectados. A criação de rotinas segue o modelo de regras gatilho-ação
típico do desenvolvimento pelo usuário final e é feita, principalmente, por
meio do aplicativo complementar instalado no
smartphone~\cite{barricelli2024routines,barricelli2022multimodal}. As
pesquisas sobre o tema concentram-se em facilitar essa autoria, comparando
a criação pelo aplicativo, por conversa e por interação multimodal em
alto-falantes com tela~\cite{barricelli2024routines,barricelli2022multimodal}.

Há, porém, uma diferença entre quem cria uma rotina e quem a vivencia. A
rotina é configurada por uma pessoa, no aplicativo associado à sua conta, e
executada em um ambiente compartilhado, onde afeta todos que estão
presentes. Estudos sobre casas inteligentes com múltiplos usuários mostram
que a instalação e a configuração tendem a se concentrar em um morador, o
que gera assimetrias de controle e tensões com os demais
residentes~\cite{geeng2019control,zeng2019multiuser}. Pessoas que trabalham
ou convivem no espaço doméstico, como babás, também são afetadas por esses
sistemas sem participar das decisões sobre eles~\cite{bernd2020bystanders}.
Sob a perspectiva da integridade contextual~\cite{nissenbaum2004privacy},
quando um dado ligado a uma pessoa, como seu horário ou sua voz, dispara
uma ação em um espaço coletivo, a informação passa a circular em um
contexto diferente daquele em que foi fornecida.

Os trabalhos citados investigam essas questões empiricamente, com
moradores, ou avaliam a usabilidade da autoria de rotinas considerando um
único autor. Recebe menos atenção a forma como a própria interface comunica
o que é uma rotina, quem a controla e a quem ela afeta. A Engenharia
Semiótica oferece uma lente para essa análise ao tratar a interface como
uma mensagem do designer para o usuário, isto é, como
metacomunicação~\cite{desouza2005semiotic}. O Método de Inspeção Semiótica
(MIS), derivado dessa teoria, permite reconstruir essa mensagem a partir
dos signos da interface, sem a participação de
usuários~\cite{desouza2006sim,desouza2009methods}, e já foi adaptado a
diferentes contextos de aplicação~\cite{fortunato2025ihc,fortunato2025adaptations}. Até onde
pudemos verificar, o método ainda não foi aplicado à criação e ao
gerenciamento de rotinas em assistentes de voz. Além disso, a MIS pressupõe
um destinatário da metamensagem, enquanto no caso das rotinas há ao menos
dois: quem cria a rotina e quem é afetado por ela.

Essa lacuna se relaciona com os Grandes Desafios de Pesquisa em IHC no
Brasil para 2025-2035 (GranDIHC-BR)~\cite{pereira2024grandihc}. O trabalho
se insere principalmente no GC4, que trata dos aspectos socioculturais em
IHC~\cite{neris2024gc4}, ao examinar uma tecnologia concebida em torno de
uma conta individual e usada em arranjos domésticos coletivos. Contribui
também para o GC1, voltado a novas abordagens teóricas e
metodológicas~\cite{dasilvajunior2024gc1}, ao propor uma adaptação da MIS
que considera múltiplos interlocutores. Por fim, dialoga com o GC7, que
aborda a interação com tecnologias emergentes como um ecossistema que
integra pessoas, tecnologias e contextos~\cite{zaina2024gc7}.

Diante disso, este artigo tem como objetivo avaliar a metacomunicação do
aplicativo Amazon Alexa na criação e no gerenciamento de rotinas,
considerando a casa como usuário coletivo, e derivar implicações de design
para lares compartilhados. A pesquisa é orientada pela seguinte questão
principal: como o aplicativo Amazon Alexa comunica a criação e o
funcionamento de rotinas, e a quem essa comunicação se dirige em uma casa
compartilhada? Dela derivam duas questões específicas:

\begin{itemize}
  \item \textbf{QP1:} Que metamensagem o designer transmite sobre autoria,
        uso de dados e pessoas afetadas pelas rotinas?
  \item \textbf{QP2:} Que potenciais rupturas de comunicação a inspeção
        identifica para a persona que administra a conta e para a persona
        que convive com as rotinas sem administrá-las?
\end{itemize}

Para responder a essas questões, aplicamos a MIS ao aplicativo Amazon
Alexa em português. A inspeção foi guiada por duas personas construídas
a partir da literatura sobre casas compartilhadas, a pessoa que
administra a conta e cria as rotinas e um morador que convive com elas
sem administrá-las, e por cenários de uso para cada uma. Em cada etapa,
os signos foram analisados quanto ao que comunicam a cada persona.
Nenhum participante foi recrutado: as personas são fictícias e a
inspeção usou contas dos próprios pesquisadores.

As contribuições deste trabalho são:

\begin{itemize}
  \item a reconstrução da metamensagem do aplicativo de Rotinas da Alexa,
        evidenciando a quem o design se dirige;
  \item um catálogo de potenciais rupturas de comunicação para quem
        administra a conta e para quem convive com as rotinas;
  \item uma adaptação da MIS que inclui uma lente para múltiplos
        interlocutores, aplicável a outros sistemas de uso coletivo;
  \item duas personas fundamentadas na literatura, com a origem de cada
        característica documentada, que podem ser reutilizadas em outros
        estudos sobre casas compartilhadas;
  \item implicações de design para rotinas em casas compartilhadas.
\end{itemize}

O restante do artigo está organizado da seguinte forma. A
Seção~\ref{sec:fundamentacao} apresenta a fundamentação teórica e os
trabalhos relacionados. A Seção~\ref{sec:metodo} descreve o método. A
Seção~\ref{sec:resultados} apresenta os resultados da inspeção,
discutidos na Seção~\ref{sec:discussao}. A Seção~\ref{sec:limitacoes}
apresenta as limitações, e a Seção~\ref{sec:conclusao}, as conclusões e
os trabalhos futuros.
